
\subsection{The involution curve and the sixty explicit core paths}

\begin{result}{B.20}{Proved}{A geometric realization of the bipyramid--pyramid incidence}
Let $U=\mathbb A^1_{\mathbb C}\setminus\{2,-2\}$ and
\[
 q_\mu(z)=z^4-\mu z^2+1.
\]
The roots of $q_\mu$, together with the distinguished point
$\infty$, define a family of five-pointed spheres with the
involution $z\mapsto-z$.  After adjoining an ordering of the
four roots, it has the explicit comparison to $M_{0,5}$ given
by (B.28).  The parameter $\mu=0$ belongs to this family:
its fibre is the square-pyramidal configuration.  Thus the
extra deletion of $0$ in cote 148, scan 39 (folio 38), does not
follow from the moduli problem with the fourth-root frame
specified on scans 37--38 (folios 36--37).

There is the following explicit realization in $\Pi_1M_{0,5}$
of the sixty $f_{B,s}$ incidences of scan 31.  In one labelling,
put
\[
 K(r)=\bigl(\infty,-ir,ir,-1,1\bigr),
 \qquad \frac1{\sqrt3}\leq r\leq1.                 \tag{B.32}
\]
Its initial configuration is bipyramidal and its final
configuration is square-pyramidal.  Relabelling (B.32) gives
sixty paths, with twenty initial objects, thirty final objects,
three paths leaving each initial object and two arriving at
each final object.  The endpoint functors are exactly
\[
 \mathfrak S_5/H_B\ \longleftarrow\
 \mathfrak S_5/H\ \longrightarrow\ \mathfrak S_5/H_Q,
 \qquad
 H_B=\langle(012),(12)(34)\rangle,\quad
 H_Q=\langle(1324)\rangle,\quad
 H=\langle(12)(34)\rangle.                         \tag{B.33}
\]
In particular these are the endpoint maps used in B.13,
now furnished with actual paths.  Formula (B.32) specifies
one realization of this part of the manuscript diagram;
it does not specify the remaining $1560$ arrows or assert
an unrecorded homotopy convention for them.
\end{result}
\begin{proof}
An involution of a complex projective line has two fixed
points.  Choose the marked one as $\infty$ and the other as
$0$.  The involution becomes $z\mapsto-z$, and its two
unmarked orbits have squared coordinates $\xi,\eta\ne0$,
with $\xi\ne\eta$.  The fourth-power frame in the manuscript
normalizes $\xi\eta=1$.  Writing $\mu=\xi+\eta$ gives
$q_\mu$ and the single condition $\mu^2\ne4$.  Conversely
the roots of this polynomial recover both orbits and their
frame.  Its discriminant is
\[
 \operatorname{disc}_z(q_\mu)=16(\mu^2-4)^2.         \tag{B.34}
\]
Changing the frame by $\zeta\in\mu_4$ changes $\mu$ by
$\zeta^2\mu$.  In particular $\mu=0$ is a valid framed
object.  Its roots are a scalar multiple of $\mu_4$; it
has the extra order-four symmetry, but this does not make
it singular or invalidate its specified order-two symmetry.
The finite cover
\[
 \{(\mu,z_1,z_2,z_3,z_4):
     q_\mu(z_j)=0,\ z_j\ne z_k\ (j\ne k)\}
 \longrightarrow U
\]
retains all labellings.  Sending its points to the labelled
sphere $(z_1,z_2,z_3,z_4,\infty)$, or to its normalized
coordinates (B.28), defines the promised map and the induced
fundamental-groupoid functor.

For the incidence, label the standard bipyramid
\[
 (1,\zeta,\zeta^2,0,\infty),\qquad
 \zeta=\exp(2\pi i/3),
\]
by $0,1,2,3,4$ in that order, and select its equatorial
point $s=0$ (the coordinate $1$).  The transformation
$M(z)=(z+1)/(z-1)$ takes it to
$K(1/\sqrt3)$.  Along (B.32) no points collide.  At the
other endpoint the four points in the finite plane are
$\{1,i,-1,-i\}$, with their complex cyclic order; the
selected point is the apex.  The bipyramid's group acts
as $H_B$ on the five labels.  At the final endpoint the
positive quarter rotation acts by $(4\,2\,3\,1)$,
which generates the same subgroup as $(1324)$.
The common half rotation is $(12)(34)$.

That common involution preserves the whole path.  No other
permutation can identify two relabelled paths: it would
preserve both endpoints, and
$H_B\cap H_Q=H$.  Conversely every element of $H$ does
identify them, by $z\mapsto-z$.  Hence the paths are
indexed by $\mathfrak S_5/H$, with endpoints the cosets
in (B.33).  This gives their cardinalities and valencies
without choices of unidentified conjugating paths.

To see the relation with the manuscript's quartic
coordinate, rescale the finite-plane coordinate in (B.32)
so that the product of the two squared orbits is $1$.
One obtains
\[
 \mu=i(r-r^{-1}),\qquad
 \mu(1/\sqrt3)=-2i/\sqrt3,\qquad \mu(1)=0.
\]
This is precisely a segment of the family just constructed.
The program \texttt{b\_quartic.py} verifies the fractional
linear transformation, discriminant, three stabilizers,
all sixty endpoint maps, and both valencies over exact
algebraic numbers.  Thus the comparison on this subdiagram
is geometric and fully specified.  Matching all the other
named arrows and their relations remains the separate
calculation recorded in B.15.
\end{proof}
